%! Author = Profesor
%! Date = 06/06/2024


\section{Ejercicios}

\subsection{Ejercicios}

\begin{exercise}
    Se dispone de\ch{NaOH} con impurezas de \ch{Ca(OH)2}. Una muestra de \masa[g]{0.77} de esta mezcla se disuelve en agua hasta obtener \SI{100}{\cubic\centi\metre} de disolución. \SI{10}{\cubic\centi\metre} de disolución necesitan \SI{20}{\cubic\centi\metre} de HCl \concentracion[m]{0.1} para su neutralización.
    Calcula el \% en peso de NaOH en la muestra.

    \begin{solu}
%        \begin{figure}[H]
%            \centering
%            \includegraphics[scale=.5]{images/imagesUD7EquilibrioAcidoBase/UD7_Problemas/general_01}
%            %\caption{}
%            %\label{fig:}
%        \end{figure}

        Sean \(x\) los moles de \ch{NaOH} e \(y\) los moles de \ch{Ca(OH)2} presentes en los \SI{100}{\centi\metre\cubed} de disolución.

        En \SI{10}{\centi\metre\cubed} (una décima parte de la muestra) hay \(x/10\) mol de \ch{NaOH} e \(y/10\) mol de \ch{Ca(OH)2}. Los moles de \ch{HCl} empleados en la neutralización son:

        \begin{equation*}
            n(\ch{HCl}) = 0,020 \ \si{\litre} \times 0,1 \ \si{\Molar} = \num{2e-3} \ \text{mol}
        \end{equation*}

        Como cada mol de \ch{NaOH} neutraliza 1 mol de \ch{HCl} y cada mol de \ch{Ca(OH)2} neutraliza 2 mol de \ch{HCl}:

        \begin{equation*}
            \dfrac{x}{10} + 2\, \dfrac{y}{10} = 2 \times 10^{-3} \ \Rightarrow \ x + 2y = 0,02 \qquad \text{(i)}
        \end{equation*}

        Además, la masa total de la muestra es (\(M(\ch{NaOH}) = 40\), \(M(\ch{Ca(OH)2}) = 74\)):

        \begin{equation*}
            40 x + 74 y = 0,77 \qquad \text{(ii)}
        \end{equation*}

        De (i): \(x = 0,02 - 2y\). Sustituyendo en (ii):

        \begin{equation*}
            40(0,02-2y) + 74y = 0,77 \ \Rightarrow \ 0,8 - 80y + 74y = 0,77 \ \Rightarrow \ -6y = -0,03 \ \Rightarrow \ y = 0,005 \ \text{mol}
        \end{equation*}

        \begin{equation*}
            x = 0,02 - 2(0,005) = 0,01 \ \text{mol de \ch{NaOH}}
        \end{equation*}

        La masa de \ch{NaOH} en la muestra es \(m = 0,01 \times 40 = \SI{0.4}{\gram}\), por lo que el porcentaje en peso de \ch{NaOH} es:

        \begin{equation*}
            \%\, \ch{NaOH} = \dfrac{0,4}{0,77} \times 100 = 51,9\ \%
        \end{equation*}

    \end{solu}
\end{exercise}

\begin{exercise}
    La densidad de un vinagre comercial (5 \% en peso de \chemfig{CH_3COOH} es de \densidad[gml]{1.0055}. Expresa la concentración del soluto en molaridad y fracción molar. ¿Qué volumen de hidróxido sódico \concentracion[m]{0.5} se necesitará para neutralizar \volumen[ml]{250} del mencionado vinagre?
    \begin{solu}
%        \begin{figure}[H]
%            \centering
%            \includegraphics[scale=.5]{images/imagesUD7EquilibrioAcidoBase/UD7_Problemas/general_02}
%            %\caption{}
%            %\label{fig:}
%        \end{figure}

        Tomamos como base de cálculo 1 L de disolución. La masa de esta disolución es:

        \begin{equation*}
            m = 1000 \ \si{\milli\litre} \times 1,0055 \ \si{\gram\per\milli\litre} = \SI{1005.5}{\gram}
        \end{equation*}

        La masa de ácido acético (5\,\% en peso) y de agua son:

        \begin{equation*}
            m(\ch{CH3COOH}) = 0,05 \times 1005,5 = \SI{50.28}{\gram} \qquad m(\ch{H2O}) = 1005,5 - 50,28 = \SI{955.22}{\gram}
        \end{equation*}

        \textbf{Molaridad}

        Con \(M(\ch{CH3COOH}) = \SI{60}{\gram\per\mole}\):

        \begin{equation*}
            n(\ch{CH3COOH}) = \dfrac{50,28}{60} = 0,838 \ \text{mol} \ \Rightarrow \ [\ch{CH3COOH}] = \SI{0.838}{\Molar}
        \end{equation*}

        \textbf{Fracción molar}

        Con \(M(\ch{H2O}) = \SI{18}{\gram\per\mole}\):

        \begin{equation*}
            n(\ch{H2O}) = \dfrac{955,22}{18} = 53,07 \ \text{mol}
        \end{equation*}

        \begin{equation*}
            x_{\ch{CH3COOH}} = \dfrac{0,838}{0,838+53,07} = 0,0155
        \end{equation*}

        \textbf{Volumen de \ch{NaOH} 0,5 M necesario}

        La reacción de neutralización es \(\ch{CH3COOH + NaOH -> CH3COONa + H2O}\) (1:1). Los moles de ácido en \volumen[ml]{250} de vinagre son:

        \begin{equation*}
            n(\ch{CH3COOH}) = 0,838 \ \si{\Molar} \times 0,250 \ \si{\litre} = 0,2095 \ \text{mol}
        \end{equation*}

        \begin{equation*}
            V(\ch{NaOH}) = \dfrac{0,2095}{0,5} = 0,419 \ \si{\litre} = \SI{418.9}{\milli\litre}
        \end{equation*}

    \end{solu}
\end{exercise}

\begin{exercise}
    Se mezclan \volumen[ml]{75} de una disolución de \ch{HCl} \concentracion[m]{0.2} con \volumen[ml]{50} de otra solución también de \ch{HCl} pero \concentracion[m]{0.5}. Con la disolución resultante se valoran \volumen[ml]{20} de una solución de \ch{NaOH}, empleándose \volumen[ml]{25}. ¿Cuál es la molaridad de la solución de hidróxido sódico?
    \begin{solu}
%        \begin{figure}[H]
%            \centering
%            \includegraphics[scale=.5]{images/imagesUD7EquilibrioAcidoBase/UD7_Problemas/general_03}
%            %\caption{}
%            %\label{fig:}
%        \end{figure}

        Al mezclar las dos disoluciones de \ch{HCl}, los moles totales de \ch{HCl} son:

        \begin{equation*}
            n(\ch{HCl}) = 0,075 \times 0,2 + 0,050 \times 0,5 = 0,015 + 0,025 = 0,04 \ \text{mol}
        \end{equation*}

        El volumen total de la mezcla es \(75+50 = \SI{125}{\milli\litre}\), por lo que:

        \begin{equation*}
            [\ch{HCl}] = \dfrac{0,04}{0,125} = \SI{0.32}{\Molar}
        \end{equation*}

        Al valorar \volumen[ml]{20} de esta disolución (\(\ch{HCl}\) es un ácido fuerte, totalmente disociado), los moles de \ch{H3O+} (y de \ch{HCl}) son:

        \begin{equation*}
            n(\ch{HCl}) = 0,32 \times 0,020 = \num{6.4e-3} \ \text{mol}
        \end{equation*}

        Como la reacción de neutralización \(\ch{HCl + NaOH -> NaCl + H2O}\) es 1:1, estos moles deben coincidir con los moles de \ch{NaOH} empleados (\volumen[ml]{25}):

        \begin{equation*}
            [\ch{NaOH}] = \dfrac{6,4 \times 10^{-3}}{0,025} = \SI{0.256}{\Molar}
        \end{equation*}

    \end{solu}
\end{exercise}

\begin{exercise}
    El \ch{CaCO3} reacciona con una disolución de \ch{HNO3}, obteniéndose \ch{CO2}, \ch{Ca(NO3)2} y \ch{H2O}. Escribe el proceso que tiene lugar. ¿Qué volumen de \ch{CO2}, medido en condiciones normales, se formará por reacción de \volumen[ml]{60} de \ch{HNO3} \concentracion[m]{2.5} con exceso de \ch{CaCO3}? ¿Qué volumen de ácido nítrico comercial, del 64\% en peso y \densidad[gml]{1.4}  de densidad, se debe tomar para preparar los \volumen[ml]{60} \concentracion[m]{2.5} a los que se refiere el apartado anterior?
    \begin{solu}
%        \begin{figure}[H]
%            \centering
%            \includegraphics[scale=.5]{images/imagesUD7EquilibrioAcidoBase/UD7_Problemas/general_04}
%            %\caption{}
%            %\label{fig:}
%        \end{figure}

        La reacción ajustada es:

        \begin{equation*}
            \ch{CaCO3 \sld{} + 2 HNO3 \aq{} -> CO2 \gas{} + Ca(NO3)2 \aq{} + H2O \lqd{}}
        \end{equation*}

        \textbf{Volumen de \ch{CO2} en condiciones normales}

        Los moles de \ch{HNO3} son:

        \begin{equation*}
            n(\ch{HNO3}) = 0,060 \ \si{\litre} \times 2,5 \ \si{\Molar} = 0,15 \ \text{mol}
        \end{equation*}

        Como la estequiometría es 2 mol de \ch{HNO3} por cada mol de \ch{CO2} formado:

        \begin{equation*}
            n(\ch{CO2}) = \dfrac{0,15}{2} = 0,075 \ \text{mol}
        \end{equation*}

        En condiciones normales (0 ºC, 1 atm), el volumen molar es \SI{22.4}{\litre\per\mole}:

        \begin{equation*}
            V(\ch{CO2}) = 0,075 \times 22,4 = \SI{1.68}{\litre}
        \end{equation*}

        \textbf{Volumen de ácido nítrico comercial necesario}

        La masa de \ch{HNO3} puro necesaria (\(M(\ch{HNO3}) = \SI{63}{\gram\per\mole}\)) es:

        \begin{equation*}
            m(\ch{HNO3}) = 0,15 \times 63 = \SI{9.45}{\gram}
        \end{equation*}

        Como el ácido comercial tiene una riqueza del 64\,\% en peso:

        \begin{equation*}
            m_{disolución} = \dfrac{9,45}{0,64} = \SI{14.77}{\gram}
        \end{equation*}

        Y, con \(\rho = \SI{1.4}{\gram\per\milli\litre}\):

        \begin{equation*}
            V = \dfrac{14,77}{1,4} = \SI{10.55}{\milli\litre}
        \end{equation*}

    \end{solu}
\end{exercise}

\subsection{Ejercicios PAU}
\begin{exerciseEVAU}
Una disolución acuosa 0,3 M de ácido hipocloroso (HClO) tiene un pH de 3,98. Calcule:
  \begin{enumerate}[label=\letraminuscula]
      \item La concentración molar de todas las especies presentes en equilibrio.
      \item El grado de disociación del ácido hipocloroso.
      \item El valor de la constante de acidez (\Ka) del ácido hipocloroso y el valor de la constante de
basicidad (\Kb) de su base conjugada.
  \end{enumerate}
\begin{solu}

\end{solu}
\end{exerciseEVAU}

\begin{exerciseEVAU}
    En una planta de tratamiento de aguas se utiliza una disolución de hidróxido sódico (NaOH) para neutralizar
contaminantes ácidos. Para ajustar el pH antes del vertido, se añade ácido nítrico (HNO3) de concentración
conocida. Se dispone de 75 mL de una disolución de NaOH 0,20 M, a la que se añade HNO3 0,10 M.
Se supone que los volúmenes son aditivos. Calcule:
\begin{enumerate}[label=\letraminuscula]
    \item El pH cuando se han añadido 50 mL de \ch{HNO3}.
    \item Escriba la reacción de neutralización. Calcule el volumen de HNO3 necesario para alcanzar el
punto de equivalencia y el pH en dicho punto.
    \item Si en lugar de ácido nítrico se utilizara ácido acético (\ch{CH3COOH}), prediga sin cálculo
numérico, el valor de pH (pH > 7 ó pH < 7 ) de una disolución acuosa de acetato sódico (\ch{CH3COONa}). Razone el resultado.
\end{enumerate}
\begin{solu}

    \textbf{a) pH al añadir 50 mL de \ch{HNO3}}

    Los moles iniciales de \ch{NaOH} y los de \ch{HNO3} añadidos son:

    \begin{equation*}
        n(\ch{NaOH}) = 0,075 \times 0,20 = 0,015 \ \text{mol} \qquad n(\ch{HNO3}) = 0,050 \times 0,10 = 0,005 \ \text{mol}
    \end{equation*}

    Como ambos son fuertes y reaccionan 1:1, queda un exceso de base:

    \begin{equation*}
        n(\ch{OH-})_{exceso} = 0,015 - 0,005 = 0,010 \ \text{mol}
    \end{equation*}

    Con \(V_T = 75+50 = \SI{125}{\milli\litre}\):

    \begin{equation*}
        [\ch{OH-}] = \dfrac{0,010}{0,125} = \SI{0.08}{\Molar} \ \Rightarrow \ pOH = -\log(0,08) = 1,10 \ \Rightarrow \ pH = 14 - 1,10 = 12,90
    \end{equation*}

    \textbf{b) Reacción de neutralización, volumen y pH en el punto de equivalencia}

    \begin{equation*}
        \ch{NaOH \aq{} + HNO3 \aq{} -> NaNO3 \aq{} + H2O \lqd{}}
    \end{equation*}

    En el punto de equivalencia, los moles de \ch{HNO3} añadidos deben igualar a los moles iniciales de \ch{NaOH} (0,015 mol):

    \begin{equation*}
        V(\ch{HNO3}) = \dfrac{0,015}{0,10} = 0,15 \ \si{\litre} = \SI{150}{\milli\litre}
    \end{equation*}

    Como se trata de la neutralización de un ácido fuerte con una base fuerte, la sal formada (\ch{NaNO3}) no se hidroliza, por lo que el \textbf{pH en el punto de equivalencia es 7} (neutro).

    \textbf{c) pH con ácido acético en lugar de ácido nítrico}

    Si se empleara ácido acético (ácido débil) en lugar de ácido nítrico, la sal formada en la neutralización sería acetato sódico (\ch{CH3COONa}). El ion acetato, \ch{CH3COO-}, es la base conjugada de un ácido débil, por lo que se hidroliza en agua:

    \begin{equation*}
        \ch{CH3COO- + H2O <=> CH3COOH + OH-}
    \end{equation*}

    Esta hidrólisis genera un exceso de \ch{OH-}, por lo que la disolución acuosa de acetato sódico tendrá \( pH > 7 \) (básico).

\end{solu}
\end{exerciseEVAU}

\begin{exerciseEVAU}
En una pequeña explotación agrícola, el responsable del invernadero utiliza una disolución de amoniaco
(\ch{NH3}) para regular el pH del agua de riego. Para preparar una nueva disolución, disuelve 7 g de amoniaco
en agua hasta obtener 500 mL de disolución.
\begin{enumerate}[label=\letraminuscula]
    \item Indique la reacción de equilibrio básico del NH3 en agua. Calcule la concentración de todas
las especies en el equilibrio y el grado de disociación del amoniaco en la disolución preparada.
    \item Determine el pH de la disolución resultante.
    \item Sin realizar cálculos numéricos, prediga si una disolución acuosa de cloruro amónico (\ch{NH4Cl})
tendrá pH ácido, básico o neutro. Justifique la respuesta explicando los procesos químicos que tienen
lugar.
\end{enumerate}
\textbf{Datos:} \Kb (\ch{NH3}) = \num{1.8e-5}
\begin{solu}

    \textbf{a) Equilibrio, concentraciones y grado de disociación}

    \begin{equation*}
        \ch{NH3 \aq{} + H2O \lqd{} <=> NH4+ \aq{} + OH- \aq{}}
    \end{equation*}

    La concentración inicial de amoniaco es (\(M(\ch{NH3}) = \SI{17}{\gram\per\mole}\)):

    \begin{equation*}
        n(\ch{NH3}) = \dfrac{7}{17} = 0,4118 \ \text{mol} \ \Rightarrow \ C_0 = \dfrac{0,4118}{0,500} = \SI{0.8235}{\Molar}
    \end{equation*}

    Llamando \(x = [\ch{NH4+}] = [\ch{OH-}]\) en el equilibrio:

    \begin{equation*}
        K_b = \dfrac{x^2}{C_0 - x} \approx \dfrac{x^2}{C_0} \ \Rightarrow \ x = \sqrt{K_b\, C_0} = \sqrt{1,8 \times 10^{-5} \times 0,8235} = \num{3.85e-3} \ \si{\Molar}
    \end{equation*}

    Por tanto:

    \begin{equation*}
        [\ch{NH4+}] = [\ch{OH-}] = \num{3.85e-3} \ \si{\Molar} \qquad [\ch{NH3}] = 0,8235 - 0,00385 = \SI{0.8197}{\Molar}
    \end{equation*}

    \begin{equation*}
        \alpha = \dfrac{x}{C_0} = \dfrac{3,85 \times 10^{-3}}{0,8235} = \num{4.68e-3} \ (0,47\ \%)
    \end{equation*}

    \textbf{b) pH de la disolución}

    \begin{equation*}
        pOH = -\log(3,85 \times 10^{-3}) = 2,41 \  \Rightarrow \  pH = 14 - 2,41 = 11,59
    \end{equation*}

    \textbf{c) pH de una disolución de \ch{NH4Cl}}

    El ion \ch{NH4+} es el ácido conjugado de la base débil \ch{NH3}, por lo que en agua se hidroliza cediendo un protón:

    \begin{equation*}
        \ch{NH4+ + H2O <=> NH3 + H3O+}
    \end{equation*}

    Esta reacción genera un exceso de \ch{H3O+}, por lo que la disolución acuosa de \ch{NH4Cl} tendrá \textbf{pH ácido} (pH \(<7\)).

\end{solu}
\end{exerciseEVAU}

\begin{exerciseEVAU}

Se dispone de 70 mL de una disolución acuosa de KOH 0,3 M, a la que se añade una disolución acuosa
de HCl 0,15 M. Suponiendo que los volúmenes son aditivos. Calcule:
    \begin{enumerate}[label=\letraminuscula]
        \item El pH cuando se han añadido 50 mL de la disolución acuosa de HCl.
        \item El volumen de la disolución acuosa de HCl que es necesario añadir a la disolución inicial de KOH para neutralizarla.
        \item Justifique el valor final del pH correspondiente al apartado b.
    \end{enumerate}

    \begin{solu}

        \textbf{a) pH al añadir 50 mL de HCl}

        \begin{equation*}
            n(\ch{KOH}) = 0,070 \times 0,3 = 0,021 \ \text{mol} \qquad n(\ch{HCl}) = 0,050 \times 0,15 = 0,0075 \ \text{mol}
        \end{equation*}

        Como ambos son fuertes (reacción 1:1), queda un exceso de base:

        \begin{equation*}
            n(\ch{OH-})_{exceso} = 0,021 - 0,0075 = 0,0135 \ \text{mol}
        \end{equation*}

        Con \(V_T = 70+50 = \SI{120}{\milli\litre}\):

        \begin{equation*}
            [\ch{OH-}] = \dfrac{0,0135}{0,120} = \SI{0.1125}{\Molar} \ \Rightarrow \ pOH = 0,95 \ \Rightarrow \ pH = 14 - 0,95 = 13,05
        \end{equation*}

        \textbf{b) Volumen de HCl para neutralizar}

        \begin{equation*}
            V(\ch{HCl}) = \dfrac{n(\ch{KOH})}{[\ch{HCl}]} = \dfrac{0,021}{0,15} = 0,14 \ \si{\litre} = \SI{140}{\milli\litre}
        \end{equation*}

        \textbf{c) pH en el punto de neutralización}

        La reacción es \(\ch{KOH + HCl -> KCl + H2O}\). Como tanto el \ch{KOH} como el \ch{HCl} son un ácido y una base fuertes, la sal formada, \ch{KCl}, procede de un ácido fuerte y una base fuerte, por lo que ni el ion \ch{K+} ni el ion \ch{Cl-} se hidrolizan en agua. Por tanto, el \textbf{pH en el punto de neutralización es 7} (disolución neutra).

    \end{solu}
\end{exerciseEVAU}

\begin{exerciseEVAU}

Mariana es una investigadora en un centro de biotecnología ambiental y está desarrollando un
método para el tratamiento de aguas contaminadas con compuestos nitrogenados provenientes de
desechos industriales. Durante un muestreo, detecta la presencia de metilamina (\ch{CH3-NH2}) en una
planta de tratamiento de efluentes, lo que podría alterar la eficiencia de los procesos biológicos
encargados de la depuración del agua. En dicho muestreo, también ha encontrado una concentración
de iones hidronio en la muestra de \( \conc{H3O+} = \qty{6.67e-12}{\M} \).


Para evaluar el impacto de la metilamina en el sistema, Mariana necesita resolver una serie de
cuestiones. ¿Podría ayudarla a resolver las siguientes?:
\begin{enumerate}[label=\letraminuscula]
    \item Escribir la reacción ácido-base de la metilamina teniendo en cuenta el valor de la
concentración de iones hidronio obtenidos en el muestreo, y calcular el pH de la disolución.
    \item Calcular la concentración inicial de metilamina.
    \item Calcular la concentración de una disolución de NaOH que tenga el mismo pH
que la disolución de metilamina.
    \item Proponer una estrategia para neutralizar la metilamina y así minimizar su
impacto en el tratamiento del agua.

Datos: \Kb (\ch{CH3NH2}) = \num{4.38e-4}


\end{enumerate}

\begin{solu}

    \textbf{a) Reacción ácido-base y pH}

    La metilamina, como el amoniaco, es una base débil que reacciona con el agua:

    \begin{equation*}
        \ch{CH3NH2 + H2O <=> CH3NH3+ + OH-}
    \end{equation*}

    Con \([\ch{H3O+}] = \num{6.67e-12} \ \si{\Molar}\):

    \begin{equation*}
        pH = -\log(6,67 \times 10^{-12}) = 11,18
    \end{equation*}

    \textbf{b) Concentración inicial de metilamina}

    \begin{equation*}
        [\ch{OH-}] = \dfrac{K_w}{[\ch{H3O+}]} = \dfrac{10^{-14}}{6,67 \times 10^{-12}} = \num{1.50e-3} \ \si{\Molar}
    \end{equation*}

    Llamando \(x = [\ch{OH-}] = [\ch{CH3NH3+}]\), y despejando \(C_0\) de la constante de basicidad sin despreciar \(x\) frente a \(C_0\) (ya que \(K_b\) no es muy pequeña):

    \begin{equation*}
        K_b = \dfrac{x^2}{C_0-x} \ \Rightarrow \ C_0 = x + \dfrac{x^2}{K_b} = 1,50\times10^{-3} + \dfrac{(1,50\times10^{-3})^2}{4,38\times10^{-4}} = \num{6.63e-3} \ \si{\Molar}
    \end{equation*}

    \textbf{c) Concentración de NaOH con el mismo pH}

    El \ch{NaOH} es una base fuerte, totalmente disociada, por lo que su concentración coincide directamente con \([\ch{OH-}]\):

    \begin{equation*}
        [\ch{NaOH}] = [\ch{OH-}] = \num{1.50e-3} \ \si{\Molar}
    \end{equation*}

    \textbf{d) Estrategia de neutralización}

    Para neutralizar la metilamina (base débil) y minimizar su impacto, se puede añadir una cantidad controlada de un ácido fuerte (por ejemplo, \ch{HCl}) en proporción estequiométrica, de modo que reaccione con la \ch{CH3NH2} formando su sal soluble, el cloruro de metilamonio (\ch{CH3NH3Cl}):

    \begin{equation*}
        \ch{CH3NH2 + HCl -> CH3NH3Cl}
    \end{equation*}

    De esta forma se consigue ajustar el pH del efluente hacia la neutralidad antes de su vertido, evitando así que la metilamina, en su forma básica libre, interfiera en los procesos biológicos de depuración.

\end{solu}

\end{exerciseEVAU}
%  AÑO 2024
\begin{exerciseEVAU}
    Se preparó una disolución de ácido acético (\ch{CH3-COOH}) disolviendo \SI{2}{\g} de éste en \SI{250}{\milli\litre}  250 mL de agua.
    Sabiendo que su constante de acidez es \( \Ka = \num{1.8e-5} \).
    \begin{enumerate}[label=\letraminuscula]
        \item Calcule la concentración de todas las especies en equilibrio.
        \item Calcule el grado de disociación del ácido acético.
        \item Calcule el pH de la disolución.
    \end{enumerate}

    \datos{} Masas atómicas relativas: C = 12; O = 16; H = 1
    \begin{solu}

        La concentración inicial de ácido acético es (\(M(\ch{CH3COOH}) = \SI{60}{\gram\per\mole}\)):

        \begin{equation*}
            n(\ch{CH3COOH}) = \dfrac{2}{60} = 0,0333 \ \text{mol} \ \Rightarrow \ C_0 = \dfrac{0,0333}{0,250} = \SI{0.1333}{\Molar}
        \end{equation*}

        \textbf{a) Concentraciones en el equilibrio}

        \begin{equation*}
            \ch{CH3COOH + H2O <=> CH3COO- + H3O+}
        \end{equation*}

        Llamando \(x = [\ch{H3O+}] = [\ch{CH3COO-}]\):

        \begin{equation*}
            K_a = \dfrac{x^2}{C_0-x} = \num{1.8e-5}
        \end{equation*}

        Resolviendo la ecuación de segundo grado (\(x^2 + 1,8\times10^{-5}x - 2,4\times10^{-6}=0\)):

        \begin{equation*}
            x = \num{1.54e-3} \ \si{\Molar}
        \end{equation*}

        \begin{equation*}
            [\ch{H3O+}] = [\ch{CH3COO-}] = \num{1.54e-3} \ \si{\Molar} \qquad [\ch{CH3COOH}] = 0,1333-0,00154 = \SI{0.1318}{\Molar}
        \end{equation*}

        \textbf{b) Grado de disociación}

        \begin{equation*}
            \alpha = \dfrac{x}{C_0} = \dfrac{1,54\times10^{-3}}{0,1333} = 0,01155 \ (1,16\ \%)
        \end{equation*}

        \textbf{c) pH de la disolución}

        \begin{equation*}
            pH = -\log(1,54 \times 10^{-3}) = 2,81
        \end{equation*}

    \end{solu}
\end{exerciseEVAU}

\begin{exerciseEVAU}
    Se añaden \SI{7}{\g}  de amoniaco a la cantidad necesaria de agua para obtener \SI{500}{\milli\litre}  500 mL de disolución.
    \begin{enumerate}[label=\letraminuscula]
        \item Calcule el grado de disociación del amoniaco.
        \item Calcule el pH de la disolución resultante.
        \item Prediga que tipo de pH (sin cálculo numérico) presentará una disolución acuosa de cloruro amónico (\ch{NH4Cl} ) y justifique la respuesta explicando los procesos químicos que tienen lugar.
    \end{enumerate}

    Datos: Masas atómicas (u): N = 14, H = 1. \Kb (\ch{NH3}) = \num{1.8e-5}
    \begin{solu}

        \begin{equation*}
            \ch{NH3 + H2O <=> NH4+ + OH-}
        \end{equation*}

        La concentración inicial es (\(M(\ch{NH3}) = \SI{17}{\gram\per\mole}\)):

        \begin{equation*}
            n(\ch{NH3}) = \dfrac{7}{17} = 0,4118 \ \text{mol} \ \Rightarrow \ C_0 = \dfrac{0,4118}{0,500} = \SI{0.8235}{\Molar}
        \end{equation*}

        \textbf{a) Grado de disociación}

        \begin{equation*}
            x = \sqrt{K_b\, C_0} = \sqrt{1,8\times10^{-5} \times 0,8235} = \num{3.85e-3} \ \si{\Molar}
        \end{equation*}

        \begin{equation*}
            \alpha = \dfrac{x}{C_0} = \dfrac{3,85\times10^{-3}}{0,8235} = 0,00468 \ (0,47\ \%)
        \end{equation*}

        \textbf{b) pH de la disolución}

        \begin{equation*}
            pOH = -\log(3,85\times10^{-3}) = 2,41 \ \Rightarrow \ pH = 14-2,41 = 11,59
        \end{equation*}

        \textbf{c) pH de una disolución de \ch{NH4Cl}}

        El ion \ch{NH4+}, ácido conjugado de la base débil \ch{NH3}, se hidroliza en agua:

        \begin{equation*}
            \ch{NH4+ + H2O <=> NH3 + H3O+}
        \end{equation*}

        Esta hidrólisis genera un exceso de \ch{H3O+}, por lo que la disolución de \ch{NH4Cl} tendrá \textbf{pH ácido}.

    \end{solu}
\end{exerciseEVAU}

\begin{exerciseEVAU}
     Una disolución de cianuro de hidrógeno (\ch{HCN}) \qty{0.01}{\M} tiene un \pH = \num{5.6}. Basándose en la reacción química correspondiente, calcule:
    \begin{enumerate}[label=\letraminuscula]
        \item El grado de disociación del \ch{HCN}.
        \item El valor de la constante de acidez.
        \item La concentración molar de todas las especies químicas presentes en el equilibrio
    \end{enumerate}
    \begin{solu}

        \begin{equation*}
            \ch{HCN + H2O <=> CN- + H3O+}
        \end{equation*}

        A partir del pH: \([\ch{H3O+}] = 10^{-5,6} = \num{2.51e-6} \ \si{\Molar}\).

        \textbf{a) Grado de disociación}

        \begin{equation*}
            \alpha = \dfrac{[\ch{H3O+}]}{C_0} = \dfrac{2,51\times10^{-6}}{0,01} = \num{2.51e-4} \ (0,0251\ \%)
        \end{equation*}

        \textbf{b) Constante de acidez}

        Como \(\alpha\) es muy pequeño, \([\ch{HCN}]_{eq} \approx C_0 = 0,01 \ \si{\Molar}\):

        \begin{equation*}
            K_a = \dfrac{[\ch{H3O+}][\ch{CN-}]}{[\ch{HCN}]} = \dfrac{(2,51\times10^{-6})^2}{0,01} = \num{6.31e-10}
        \end{equation*}

        \textbf{c) Concentraciones en el equilibrio}

        \begin{equation*}
            [\ch{H3O+}] = [\ch{CN-}] = \num{2.51e-6} \ \si{\Molar} \qquad [\ch{HCN}] \approx \SI{0.01}{\Molar}
        \end{equation*}

        \begin{equation*}
            [\ch{OH-}] = \dfrac{K_w}{[\ch{H3O+}]} = \dfrac{10^{-14}}{2,51\times10^{-6}} = \num{3.98e-9} \ \si{\Molar}
        \end{equation*}

    \end{solu}
\end{exerciseEVAU}

\begin{exerciseEVAU}

    Dadas las siguientes especies químicas: \ch{CH3COOH}, \ch{HCO3\mch[]}, \ch{NH3} y \ch{HNO3}, justifique, según
    la teoría de Brönsted-Lowry:
    \begin{enumerate}[label=\letraminuscula]
        \item Cuáles pueden actuar sólo como ácidos y cuáles sólo como bases.
        \item Cuáles pueden actuar como ácidos y como bases (anfótero)
    \end{enumerate}
    \begin{solu}

        Según la teoría de Brönsted-Lowry, un ácido es toda especie capaz de ceder un protón (\ch{H+}), y una base es toda especie capaz de aceptarlo.

        \begin{itemize}
            \item \(\ch{CH3COOH}\): puede ceder un protón (\(\ch{CH3COOH -> CH3COO- + H+}\)), pero no dispone de pares de electrones no compartidos accesibles ni de hidrógenos adicionales para aceptar un protón en condiciones habituales. \textbf{Actúa solo como ácido.}
            \item \(\ch{HNO3}\): ácido fuerte, cede fácilmente su protón (\(\ch{HNO3 -> NO3- + H+}\)) y no puede aceptar otro. \textbf{Actúa solo como ácido.}
            \item \(\ch{NH3}\): dispone de un par de electrones no compartido en el nitrógeno que le permite aceptar un protón (\(\ch{NH3 + H+ -> NH4+}\)), pero no tiene ningún hidrógeno ácido que ceder en disolución acuosa. \textbf{Actúa solo como base.}
        \end{itemize}

        \textbf{b) Especie anfótera}

        \begin{itemize}
            \item \(\ch{HCO3-}\): puede ceder un protón, actuando como ácido (\(\ch{HCO3- -> CO3^2- + H+}\)), o aceptar un protón, actuando como base (\(\ch{HCO3- + H+ -> H2CO3}\)). Por tanto, es una especie \textbf{anfótera}.
        \end{itemize}

    \end{solu}
\end{exerciseEVAU}

\begin{exerciseEVAU}
    Disponemos de dos disoluciones acuosas, una disolución de \concentracion[m]{0.10} de \ch{NaOH} y otra y otra disolución \concentracion[m]{0.5} de \ch{ HCl}
    \begin{enumerate}[label=\letraminuscula]
        \item Calcule el pH de cada una de las disoluciones.
        \item Calcular el pH de la disolución resultante si mezclamos 80 mL de la primera disolución y 20 mL de la segunda disolución, considere volúmenes aditivos.
        \item Calcular el volumen de disolución de \ch{HCl} que hay que añadir a 100 mL de la disolución de NaOH para neutralizarla.
    \end{enumerate}
    \solutionCM{a) pH (\ch{HCl}) = \num{0.3} y pH (\ch{NaOH}) = \num{13}; b) pH = \num{1.7}; V = \volumen[l]{0.02}}

    \begin{solu}
        Se la reacción:
        \begin{reaction*}
            NaOH + HCl -> NaCl + H2O
        \end{reaction*}
        En este caso que tenemos que tanto el \ch{NaOH} como el \ch{HCl} son fuertes, por lo tanto están completamente disociados. Podemos decir que la concentración de iones \ch{OH\mch[]} y \ch{H3O\pch[]} son iguales a las concentraciones iniciales de \ch{NaOH} y \ch{HCl} respectivamente.
        \begin{reactions*}
            NaOH & -> OH\mch[] + Na\\
            HCl + H2O & -> H3O\pch[] + Cl-\\
        \end{reactions*}
        En este caso tendremos que:
        % \begin{itemize}
        %     \item \(\conc{NaOH}=\conc{OH-} = 0.1\)

        %     \( pOH = -\log{0.1} = 1 \rightarrow pOH + pH = 14 \rightarrow pH = 14 - pOH = 14 - 1 \rightarrow pH = 13\).
        %     \item \(\conc{HCl}=\conc{H3O+} = 0.5\)

        %     \( pH = -\log{0.5} = 0.3\).
        % \end{itemize}
%        \begin{figure}[H]
%            \centering
%            \includegraphics[scale=.5]{images/imagesUD7EquilibrioAcidoBase/UD7_Problemas/EvAU_01}
%            %\caption{}
%            %\label{fig:}
%        \end{figure}

        \begin{enumerate}[label=\letraminuscula]
            \item El \ch{NaOH} es una base fuerte y el \ch{HCl} es un ácido fuerte, por lo que ambos están totalmente disociados. Por ello, las concentraciones de iones hidroxilo e iones hidronio son iguales a las de \ch{NaOH} y \ch{HCl}:
            \begin{align*}
                \conc{OH-}  &= \num{0.1}\,\unit{\M} & \mathrm{pOH} &= -\log \num{0.1} = 1 & & \textcolor{red}{\mathrm{pH} = 13} \\
                \conc{H3O+} &= \num{0.5}\,\unit{\M} &              &                      & & \textcolor{red}{\mathrm{pH} = -\log \num{0.5} = \num{0.3}}
            \end{align*}

            \item Moles iniciales de cada especie:
            \begin{align*}
                n_0(\ch{OH-})  &= M \cdot V = \num{0.1} \cdot \num{0.08} = \qty{0.008}{\mol} \\
                n_0(\ch{H3O+}) &= M \cdot V = \num{0.5} \cdot \num{0.02} = \qty{0.01}{\mol}
            \end{align*}
            Después de la neutralización quedan \(\num{0.01} - \num{0.008} = \qty{0.002}{\mol}\) de iones hidronio en \(\qty{80}{\mL} + \qty{20}{\mL} = \qty{0.1}{\litre}\) de disolución:
            \[
                \conc{H3O+} = \frac{\qty{0.002}{\mol}}{\qty{0.1}{\litre}} = \qty{0.02}{\M} \qquad \textcolor{red}{\mathrm{pH} = -\log \num{0.02} = \num{1.7}}
            \]

            \item Para neutralizar, los moles de \ch{HCl} deben ser iguales a los de \ch{NaOH}:
            \begin{gather*}
                n(\ch{NaOH}) = M \cdot V = \num{0.1} \cdot \num{0.1} = \qty{0.01}{\mol} \\
                n(\ch{HCl}) = n(\ch{NaOH}) = \qty{0.01}{\mol} = M \cdot V \\
                V = \frac{\num{0.01}}{\num{0.5}} = \qty{0.02}{\litre} = \textcolor{red}{\qty{20}{\mL}}
            \end{gather*}
        \end{enumerate}

    \end{solu}

\end{exerciseEVAU}

\begin{exerciseEVAU}
    Se disuelven \masa[g]{23} de ácido metanoico, \ch{HCOOH}, en agua hasta obtener 10 litros de disolución.
    La concentración de iones \ch{H3O+} es \num{0,003} M. Calcule:

    \begin{enumerate}[label=\letraminuscula]
        \item El pH de la disolución y el grado de disociación.
        \item La constante \Ka{} del ácido metanoico.
    \end{enumerate}

    Masas atómicas: H = 1; C = 12; O = 16.

    \solutionCM{a) \(\alpha = 0.06\), pH = 2.52; b) \(K_a = \num{1.9e-4}\)}
    \begin{solu}
%        \begin{figure}[H]
%            \centering
%            \includegraphics[scale=.5]{images/imagesUD7EquilibrioAcidoBase/UD7_Problemas/EvAU_02.png}
%            %\caption{}
%            %\label{fig:}
%        \end{figure}

        La concentración inicial de ácido metanoico es (\(M(\ch{HCOOH}) = \SI{46}{\gram\per\mole}\)):

        \begin{equation*}
            n(\ch{HCOOH}) = \dfrac{23}{46} = 0,5 \ \text{mol} \ \Rightarrow \ C_0 = \dfrac{0,5}{10} = \SI{0.05}{\Molar}
        \end{equation*}

        \textbf{a) pH y grado de disociación}

        \begin{equation*}
            pH = -\log(0,003) = 2,52
        \end{equation*}

        \begin{equation*}
            \alpha = \dfrac{[\ch{H3O+}]}{C_0} = \dfrac{0,003}{0,05} = 0,06 \ (6\ \%)
        \end{equation*}

        \textbf{b) Constante \(K_a\)}

        \begin{equation*}
            \ch{HCOOH + H2O <=> HCOO- + H3O+}
        \end{equation*}

        \begin{equation*}
            K_a = \dfrac{[\ch{H3O+}][\ch{HCOO-}]}{[\ch{HCOOH}]} = \dfrac{(0,003)^2}{0,05-0,003} = \dfrac{9\times10^{-6}}{0,047} = \num{1.9e-4}
        \end{equation*}

    \end{solu}
\end{exerciseEVAU}


\begin{exerciseEVAU}
    Se dispone de dos disoluciones, una de \ch{HCl} \concentracion[m]{0,1}  y otra de \ch{NaOH} \concentracion[m]{0.05} a) Calcule el pH de cada una de ellas. b) ¿Qué pH tendrá la mezcla de \volumen[ml]{500} de cada una de las disoluciones? c) Calcule el volumen de disolución de  \ch{NaOH} que hay que añadir a \volumen[ml]{100} de la de \ch{HCl} para neutralizarla. En todos los casos suponga volúmenes aditivos.

    \solutionCM{a) \pH (\ch{HCl}) = 1, \pH (\ch{NaOH}) = 12.70; b) \pH = 1.60; c) V\subin{\ch{NaOH}} = \volumen[ml]{200}}

    \begin{solu}

%        \begin{figure}[H]
%            \centering
%            \includegraphics[scale=.5]{images/imagesUD7EquilibrioAcidoBase/UD7_Problemas/EvAU_03.png}
%            %\caption{}
%            %\label{fig:}
%        \end{figure}

        \textbf{a) pH de cada disolución}

        \begin{equation*}
            pH(\ch{HCl}) = -\log(0,1) = 1
        \end{equation*}

        \begin{equation*}
            pOH(\ch{NaOH}) = -\log(0,05) = 1,30 \ \Rightarrow \ pH(\ch{NaOH}) = 14-1,30 = 12,70
        \end{equation*}

        \textbf{b) pH de la mezcla de 500 mL de cada una}

        \begin{equation*}
            n(\ch{HCl}) = 0,5\times0,1 = 0,05 \ \text{mol} \qquad n(\ch{NaOH}) = 0,5\times0,05 = 0,025 \ \text{mol}
        \end{equation*}

        \begin{equation*}
            n(\ch{H+})_{exceso} = 0,05-0,025 = 0,025 \ \text{mol} \ \Rightarrow \ [\ch{H+}] = \dfrac{0,025}{1,0} = \SI{0.025}{\Molar}
        \end{equation*}

        \begin{equation*}
            pH = -\log(0,025) = 1,60
        \end{equation*}

        \textbf{c) Volumen de \ch{NaOH} para neutralizar 100 mL de \ch{HCl}}

        \begin{equation*}
            n(\ch{HCl}) = 0,1\times0,1 = 0,01 \ \text{mol} \ \Rightarrow \ V(\ch{NaOH}) = \dfrac{0,01}{0,05} = 0,2 \ \si{\litre} = \SI{200}{\milli\litre}
        \end{equation*}

    \end{solu}
\end{exerciseEVAU}

%\solutionCM{} a) pH (\ch{HCl}) = \num{1.0}, pH (\ch{NaOH}) = \num{12.7}; b) pH = \num{1.6}; c) V = \volumen{ml}{200}

\begin{exerciseEVAU}
    El grado de disociación de una disolución de ácido acético en agua es del \num{2.53}\%.
    \begin{enumerate}[label=\letraminuscula]
        \item Escriba la ecuación de disociación del ácido acético.
        \item Calcule la concentración inicial de ácido acético antes de disociarse.
        \item Calcule el pH de la disolución resultante.
    \end{enumerate}
    \datos{}: \(K_a = \num{1.76e-5}\)

    \solutionCM{a) \ch{CH3COOH + H2O <=> CH3COO^{-} + H3O^{+}}; b) \conc{CH3COOH}\subin{0} = \concentracion[m]{0.027} ;c) \ph{3.17}}

    \begin{solu}

%        \begin{figure}[H]
%            \centering
%            \includegraphics[scale=.5]{images/imagesUD7EquilibrioAcidoBase/UD7_Problemas/EvAU_04.png}
%            %\caption{}
%            %\label{fig:}
%        \end{figure}

        \textbf{a) Ecuación de disociación}

        \begin{equation*}
            \ch{CH3COOH + H2O <=> CH3COO- + H3O+}
        \end{equation*}

        \textbf{b) Concentración inicial}

        A partir de la relación exacta \(K_a = \dfrac{\alpha^2\, C_0}{1-\alpha}\), con \(\alpha = 0,0253\):

        \begin{equation*}
            C_0 = \dfrac{K_a\,(1-\alpha)}{\alpha^2} = \dfrac{1,76\times10^{-5}\times(1-0,0253)}{(0,0253)^2} = \SI{0.027}{\Molar}
        \end{equation*}

        \textbf{c) pH de la disolución}

        \begin{equation*}
            [\ch{H3O+}] = \alpha\, C_0 = 0,0253 \times 0,027 = \num{6.78e-4} \ \si{\Molar}
        \end{equation*}

        \begin{equation*}
            pH = -\log(6,78\times10^{-4}) = 3,17
        \end{equation*}

    \end{solu}
\end{exerciseEVAU}



\begin{exerciseEVAU}
    Disponemos de dos disoluciones una de \ch{HNO3} \concentracion[m]{0.5} y otra de \ch{NaOH} \concentracion[m]{0.4}.
    \begin{enumerate}[label=\letraminuscula]
        \item Calcule el pH de cada una de ellas.
        \item ¿Qué pH tendría la mezcla de \volumen[ml]{100} de cada una de las disoluciones
        \item Calcule el volumen de la disolución de \ch{NaOH} \concentracion[m]{0.4} que hay que añadir a \volumen[ml]{100} de  \ch{HNO3} \concentracion[m]{0.5} para neutralizarla.
    \end{enumerate}
    Suponer en todos los casos, volúmenes aditivos.

    \solutionCM{a) pH (\ch{NaOH})= \num{13.6}, pH (\ch{HCl})= \num{0.3}; b) \ph{1.3}; c) V = \volumen[ml]{125}}

    \begin{solu}

%        \begin{figure}[H]
%            \centering
%            \includegraphics[scale=.5]{images/imagesUD7EquilibrioAcidoBase/UD7_Problemas/EvAU_05.png}
%            %\caption{}
%            %\label{fig:}
%        \end{figure}

        \textbf{a) pH de cada disolución}

        \begin{equation*}
            pH(\ch{HNO3}) = -\log(0,5) = 0,30
        \end{equation*}

        \begin{equation*}
            pOH(\ch{NaOH}) = -\log(0,4) = 0,40 \ \Rightarrow \ pH(\ch{NaOH}) = 14-0,40 = 13,60
        \end{equation*}

        \textbf{b) pH de la mezcla de 100 mL de cada una}

        \begin{equation*}
            n(\ch{HNO3}) = 0,1\times0,5 = 0,05 \ \text{mol} \qquad n(\ch{NaOH}) = 0,1\times0,4 = 0,04 \ \text{mol}
        \end{equation*}

        \begin{equation*}
            n(\ch{H+})_{exceso} = 0,05-0,04 = 0,01 \ \text{mol} \ \Rightarrow \ [\ch{H+}] = \dfrac{0,01}{0,2} = \SI{0.05}{\Molar}
        \end{equation*}

        \begin{equation*}
            pH = -\log(0,05) = 1,30
        \end{equation*}

        \textbf{c) Volumen de \ch{NaOH} para neutralizar 100 mL de \ch{HNO3}}

        \begin{equation*}
            n(\ch{HNO3}) = 0,1\times0,5 = 0,05 \ \text{mol} \ \Rightarrow \ V(\ch{NaOH}) = \dfrac{0,05}{0,4} = 0,125 \ \si{\litre} = \SI{125}{\milli\litre}
        \end{equation*}

    \end{solu}
\end{exerciseEVAU}


\begin{exerciseEVAU}
    Se preparó una disolución que contenía \masa[g]{2.48} de amoniaco en un volumen de 1 litro de agua.
    \begin{enumerate}[label=\letraminuscula]
        \item Escriba la ecuación de hidrólisis del amoniaco.
        \item Calcule el grado de disociación del amoniaco.
        \item Calcule el pH de la disolución resultante.
    \end{enumerate}
    \datos{}: \(k_b = \num{1.81e-5}\); Masas atómicas: N = \num{14.00}; H = \num{1.00}

    \solutionCM{a) \ch{NH3 + H2O <=> NH4+ + OH-}; b) \(\alpha = 1.1\%\) ;c) \(\pH = 11.21\)}

    \begin{solu}

%        \begin{figure}[H]
%            \centering
%            \includegraphics[scale=.5]{images/imagesUD7EquilibrioAcidoBase/UD7_Problemas/EvAU_06.png}
%            %\caption{}
%            %\label{fig:}
%        \end{figure}

        \textbf{a) Ecuación de hidrólisis}

        \begin{equation*}
            \ch{NH3 + H2O <=> NH4+ + OH-}
        \end{equation*}

        La concentración inicial es (\(M(\ch{NH3}) = 14,00+3(1,00) = \SI{17.00}{\gram\per\mole}\)):

        \begin{equation*}
            n(\ch{NH3}) = \dfrac{2,48}{17,00} = 0,1459 \ \text{mol} \ \Rightarrow \ C_0 = \dfrac{0,1459}{1} = \SI{0.1459}{\Molar}
        \end{equation*}

        \textbf{b) Grado de disociación}

        \begin{equation*}
            x = \sqrt{k_b\, C_0} = \sqrt{1,81\times10^{-5}\times0,1459} = \num{1.625e-3} \ \si{\Molar}
        \end{equation*}

        \begin{equation*}
            \alpha = \dfrac{x}{C_0} = \dfrac{1,625\times10^{-3}}{0,1459} = 0,0111 \ (1,1\ \%)
        \end{equation*}

        \textbf{c) pH de la disolución}

        \begin{equation*}
            pOH = -\log(1,625\times10^{-3}) = 2,79 \ \Rightarrow \ pH = 14-2,79 = 11,21
        \end{equation*}

    \end{solu}
\end{exerciseEVAU}

\begin{exerciseEVAU}
    Se prepara una disolución disolviendo \masa[g]{4} de \ch{NaOH} en \volumen[l]{2} de agua.
    \begin{enumerate}[label=\letraminuscula]
        \item Calcule el pH de la disolución.
        \item Si ahora se le añaden \volumen[ml]{500} de disolución \concentracion[m]{0.5} de \ch{HCl}, ¿cuál es el pH de la disolución resultante?.
        \item Calcule el volumen de disolución \concentracion[m]{0.5} de \ch{HCl} necesario para neutralizar \volumen[ml]{50.0} de la disolución inicial de  \ch{NaOH}.
    \end{enumerate}
    \datos{}: \ch{Na} = \num{23.0}, \ch{H} = \num{1}, \ch{O} = \num{16}, \ch{Cl} = \num{35.5}

    \solutionCM{a= \ph{12.7}; b) \ph{1.22}; c) \volumen[ml]{25}}

    \begin{solu}
%
%        \begin{figure}[H]
%            \centering
%            \includegraphics[scale=.5]{images/imagesUD7EquilibrioAcidoBase/UD7_Problemas/EvAU_07.png}
%            %\caption{}
%            %\label{fig:}
%        \end{figure}

        La concentración inicial de \ch{NaOH} es (\(M(\ch{NaOH}) = 23,0+16+1 = \SI{40.0}{\gram\per\mole}\)):

        \begin{equation*}
            n(\ch{NaOH}) = \dfrac{4}{40} = 0,1 \ \text{mol} \ \Rightarrow \ C_0 = \dfrac{0,1}{2} = \SI{0.05}{\Molar}
        \end{equation*}

        \textbf{a) pH de la disolución inicial}

        \begin{equation*}
            pOH = -\log(0,05) = 1,30 \ \Rightarrow \ pH = 14-1,30 = 12,70
        \end{equation*}

        \textbf{b) pH tras añadir 500 mL de HCl 0,5 M}

        \begin{equation*}
            n(\ch{HCl}) = 0,5\times0,5 = 0,25 \ \text{mol} \ \Rightarrow \ n(\ch{H+})_{exceso} = 0,25-0,1 = 0,15 \ \text{mol}
        \end{equation*}

        Con \(V_T = 2+0,5 = \SI{2.5}{\litre}\):

        \begin{equation*}
            [\ch{H+}] = \dfrac{0,15}{2,5} = \SI{0.06}{\Molar} \ \Rightarrow \ pH = -\log(0,06) = 1,22
        \end{equation*}

        \textbf{c) Volumen de HCl para neutralizar 50,0 mL de la disolución inicial}

        \begin{equation*}
            n(\ch{NaOH})_{50\,mL} = 0,05 \times 0,050 = \num{2.5e-3} \ \text{mol}
        \end{equation*}

        \begin{equation*}
            V(\ch{HCl}) = \dfrac{2,5\times10^{-3}}{0,5} = \num{5e-3} \ \si{\litre} = \SI{5.0}{\milli\litre}
        \end{equation*}

    \end{solu}
\end{exerciseEVAU}

\begin{exerciseEVAU}
    Se dispone de dos disoluciones: una de \ch{HNO3} \concentracion[m]{0.5} y otra de \ch{NaOH} \concentracion[m]{0.25}.
    \begin{enumerate}[label=\letraminuscula]
        \item Calcula el pH de ambas disoluciones.
        \item Se valoran \volumen[ml]{20} de la disolución de \ch{HNO3} con la disolución de \ch{NaOH}. ¿Qué volumen de ésta será necesario añadir para alcanzar el punto final de la valoración? ¿Qué reacción tiene lugar en la misma?
        \item ¿Cuánto valdrá el pH cuando se hayan añadido \volumen[ml]{20} de \ch{NaOH}
    \end{enumerate}

    \begin{solu}
%
%        \begin{figure}[H]
%            \centering
%            \includegraphics[scale=.5]{images/imagesUD7EquilibrioAcidoBase/UD7_Problemas/EvAU_08.png}
%            %\caption{}
%            %\label{fig:}
%        \end{figure}

        \textbf{a) pH de ambas disoluciones}

        \begin{equation*}
            pH(\ch{HNO3}) = -\log(0,5) = 0,30
        \end{equation*}

        \begin{equation*}
            pOH(\ch{NaOH}) = -\log(0,25) = 0,60 \ \Rightarrow \ pH(\ch{NaOH}) = 14-0,60 = 13,40
        \end{equation*}

        \textbf{b) Volumen de NaOH para alcanzar el punto de equivalencia}

        La reacción es \(\ch{HNO3 + NaOH -> NaNO3 + H2O}\). Con \(n(\ch{HNO3}) = 0,020\times0,5 = 0,01\) mol:

        \begin{equation*}
            V(\ch{NaOH}) = \dfrac{0,01}{0,25} = 0,04 \ \si{\litre} = \SI{40}{\milli\litre}
        \end{equation*}

        \textbf{c) pH tras añadir 20 mL de NaOH}

        Como aún no se ha añadido el volumen de equivalencia (40 mL), queda ácido en exceso:

        \begin{equation*}
            n(\ch{NaOH})_{añadido} = 0,020\times0,25 = 0,005 \ \text{mol} \ \Rightarrow \ n(\ch{H+})_{exceso} = 0,01-0,005 = 0,005 \ \text{mol}
        \end{equation*}

        Con \(V_T = 20+20 = \SI{40}{\milli\litre}\):

        \begin{equation*}
            [\ch{H+}] = \dfrac{0,005}{0,040} = \SI{0.125}{\Molar} \ \Rightarrow \ pH = -\log(0,125) = 0,90
        \end{equation*}

    \end{solu}
\end{exerciseEVAU}

\begin{exerciseEVAU}
    Se dispone en el laboratorio de una disolución de ácido benzoico (disolución A), de concentración desconocida. \volumen[ml]{100} de esta disolución se diluyen con agua hasta un volumen total de \volumen[ml]{500} para obtener otra disolución (disolución B) cuyo pH resulta ser \num{2.90}. Calcula: a) Las concentraciones molares de todas las especies presentes en la disolución B. b) El grado de disociación del ácido benzoico en la disolución B. c) La concentración inicial de ácido benzoico en la disolución A.

    \datos{}: \(K_a = \num{6.3e-5}\)

    \solutionCM{a) \conc{PhCOOH}\subin{eq} = \concentracion[m]{0.0252}; \(\conc{H3O+} = \conc{PhCOO-} = \concentracion[m]{1.26e-3}\)  b) \(\alpha = \num{0.051}\), despreciando el grado de disociación con respecto a la unidad; c) \concentracion[m]{0.1323}}

    \begin{solu}

%        \begin{figure}[H]
%            \centering
%            \includegraphics[scale=.5]{images/imagesUD7EquilibrioAcidoBase/UD7_Problemas/EvAU_09.png}
%            %\caption{}
%            %\label{fig:}
%        \end{figure}

        La disolución B (\SI{500}{\milli\litre}) se obtuvo diluyendo \SI{100}{\milli\litre} de la disolución A, con un factor de dilución de \(500/100 = 5\).

        \textbf{a) Concentraciones en la disolución B}

        \begin{equation*}
            [\ch{H3O+}] = 10^{-2,90} = \num{1.259e-3} \ \si{\Molar} = [\ch{PhCOO-}]
        \end{equation*}

        Llamando \(C_B\) a la concentración inicial de ácido benzoico en B:

        \begin{equation*}
            K_a = \dfrac{[\ch{H3O+}]^2}{C_B - [\ch{H3O+}]} \ \Rightarrow \ C_B = [\ch{H3O+}] + \dfrac{[\ch{H3O+}]^2}{K_a} = 1,259\times10^{-3} + \dfrac{(1,259\times10^{-3})^2}{6,3\times10^{-5}}
        \end{equation*}

        \begin{equation*}
            C_B = \num{2.64e-2} \ \si{\Molar} \ \Rightarrow \ [\ch{PhCOOH}]_{eq} = 0,02642 - 0,00126 = \num{2.52e-2} \ \si{\Molar}
        \end{equation*}

        \textbf{b) Grado de disociación en B}

        \begin{equation*}
            \alpha = \dfrac{[\ch{H3O+}]}{C_B} = \dfrac{1,259\times10^{-3}}{0,02642} = 0,0477 \ (4,8\ \%)
        \end{equation*}

        \textbf{c) Concentración inicial en la disolución A}

        Como B se obtuvo diluyendo A cinco veces, la concentración de A era 5 veces mayor:

        \begin{equation*}
            C_A = 5 \times C_B = 5 \times 0,02642 = \SI{0.132}{\Molar}
        \end{equation*}

    \end{solu}
\end{exerciseEVAU}

\begin{exerciseEVAU}
    Se prepara \volumen[ml]{100} de disolución \concentracion[m]{0.15} de \ch{HCl}.
    \begin{enumerate}[label=\letraminuscula]
        \item ¿Cuánto vale el pH de la disolución?
        \item  ¿Qué volumen de una disolución 0,25 M de NaOH será necesario para neutralizar los 100 mL del ácido? Escribe la reacción que tiene lugar.
        \item ¿Qué volumen de HCl comercial del 35 \% en masa y densidad 1,18 g/mL ha sido necesario para preparar por dilución los 100 mL de la disolución 0,15 M de partida?
    \end{enumerate}

    \datos{} Masas atómicas: Cl=35,5; H=1,0

    \solutionCM{ a) \(\pH = 0.82\); b) \(V_{\ch{NaOH}} = \volumen[ml]{60}\) ; c) \(V_{\ch{HCl}} = \volumen[ml]{1.33}\)}

    \begin{solu}

%        \begin{figure}[H]
%            \centering
%            \includegraphics[scale=.5]{images/imagesUD7EquilibrioAcidoBase/UD7_Problemas/EvAU_10.png}
%            %\caption{}
%            %\label{fig:}
%        \end{figure}

        \textbf{a) pH de la disolución}

        \begin{equation*}
            pH = -\log(0,15) = 0,82
        \end{equation*}

        \textbf{b) Volumen de NaOH necesario}

        \begin{equation*}
            \ch{HCl + NaOH -> NaCl + H2O}
        \end{equation*}

        \begin{equation*}
            n(\ch{HCl}) = 0,1\times0,15 = 0,015 \ \text{mol} \ \Rightarrow \ V(\ch{NaOH}) = \dfrac{0,015}{0,25} = 0,06 \ \si{\litre} = \SI{60}{\milli\litre}
        \end{equation*}

        \textbf{c) Volumen de HCl comercial necesario}

        La masa de \ch{HCl} puro necesaria (\(M(\ch{HCl}) = 35,5+1,0 = \SI{36.5}{\gram\per\mole}\)) es:

        \begin{equation*}
            m(\ch{HCl}) = 0,015\times36,5 = \SI{0.5475}{\gram}
        \end{equation*}

        Con una riqueza del 35\,\% en peso:

        \begin{equation*}
            m_{disolución} = \dfrac{0,5475}{0,35} = \SI{1.564}{\gram} \ \Rightarrow \ V = \dfrac{1,564}{1,18} = \SI{1.33}{\milli\litre}
        \end{equation*}

    \end{solu}
\end{exerciseEVAU}

\begin{exerciseEVAU}
    Se dispone de disoluciones de igual concentración de las siguientes sales: \ch{NaCl} y \ch{NH4Cl} ¿Cuál de ellas tendrá mayor pH? Justifica la respuesta escribiendo las correspondientes reacciones de equilibrio.

    \datos{}: \(K_b(\ch{NH3}) = \num{1.8e-5}\)

    \solutionCM{ \(\pH \ch{Nacl} > \pH \ch{NH4Cl}\) }

    \begin{solu}
%
%        \begin{figure}[H]
%            \centering
%            \includegraphics[scale=.5]{images/imagesUD7EquilibrioAcidoBase/UD7_Problemas/EvAU_11.png}
%            %\caption{}
%            %\label{fig:}
%        \end{figure}

        El \ch{NaCl} procede de un ácido fuerte (\ch{HCl}) y una base fuerte (\ch{NaOH}), por lo que ni el \ch{Na+} ni el \ch{Cl-} se hidrolizan en agua: la disolución es \textbf{neutra} (\(pH = 7\)).

        El \ch{NH4Cl} procede de una base débil (\ch{NH3}) y un ácido fuerte (\ch{HCl}). El ion \ch{NH4+} es el ácido conjugado de \ch{NH3} y se hidroliza:

        \begin{equation*}
            \ch{NH4+ + H2O <=> NH3 + H3O+}
        \end{equation*}

        Esta hidrólisis genera un exceso de \ch{H3O+}, por lo que la disolución de \ch{NH4Cl} es \textbf{ácida} (\(pH < 7\)).

        Por tanto, \( pH(\ch{NaCl}) > pH(\ch{NH4Cl}) \).

    \end{solu}
\end{exerciseEVAU}

\begin{exerciseEVAU}
    Se dispone de dos muestras de \volumen[ml]{100}  de dos disoluciones distintas, una \concentracion[m]{0.1}  en \ch{HNO3} y otra \concentracion[m]{0.5} en \ch{KOH}.
    \begin{enumerate}[label=\letraminuscula]
        \item Calcula el pH de cada disolución.
        \item ¿Qué reacción tendrá lugar al mezclar las? Escribe la correspondiente ecuación química.
        \item ¿Qué pH tendrá la disolución resultante de la mezcla? Supón volúmenes aditivos
    \end{enumerate}
    \solutionCM{a) a) pH (\ch{KOH})= \num{13.7}, a) pH (\ch{HNO3})= \num{1.0}; b) \ch{HNO3 + KOH -> KNO3 + H2O}; c) \ph{13.3}}

    \begin{solu}

%        \begin{figure}[H]
%            \centering
%            \includegraphics[scale=.5]{images/imagesUD7EquilibrioAcidoBase/UD7_Problemas/EvAU_12.png}
%            %\caption{}
%            %\label{fig:}
%        \end{figure}

        \textbf{a) pH de cada disolución}

        \begin{equation*}
            pH(\ch{HNO3}) = -\log(0,1) = 1,0
        \end{equation*}

        \begin{equation*}
            pOH(\ch{KOH}) = -\log(0,5) = 0,30 \ \Rightarrow \ pH(\ch{KOH}) = 14-0,30 = 13,70
        \end{equation*}

        \textbf{b) Reacción al mezclarlas}

        \begin{equation*}
            \ch{HNO3 + KOH -> KNO3 + H2O}
        \end{equation*}

        \textbf{c) pH de la mezcla}

        \begin{equation*}
            n(\ch{HNO3}) = 0,1\times0,1 = 0,01 \ \text{mol} \qquad n(\ch{KOH}) = 0,1\times0,5 = 0,05 \ \text{mol}
        \end{equation*}

        \begin{equation*}
            n(\ch{OH-})_{exceso} = 0,05-0,01 = 0,04 \ \text{mol} \ \Rightarrow \ [\ch{OH-}] = \dfrac{0,04}{0,2} = \SI{0.2}{\Molar}
        \end{equation*}

        \begin{equation*}
            pOH = -\log(0,2) = 0,70 \ \Rightarrow \ pH = 14-0,70 = 13,30
        \end{equation*}

    \end{solu}
\end{exerciseEVAU}

\begin{exerciseEVAU}
    Ordena de menor a mayor pH las disoluciones \concentracion[m]{0.1}  de los siguientes compuestos: a) \ch{NH4Cl}; b) \ch{NaOH}, c) \ch{NaCl}, d) \ch{NH3}. Justifica tu respuesta.

    \solutionCM{ \(\ch{NaOH} > \ch{NH3} > \ch{NaCl} > \ch{NH4Cl} \) }

    \begin{solu}

%        \begin{figure}[H]
%            \centering
%            \includegraphics[scale=.5]{images/imagesUD7EquilibrioAcidoBase/UD7_Problemas/EvAU_13.png}
%            %\caption{}
%            %\label{fig:}
%        \end{figure}

        \begin{itemize}
            \item \(\ch{NaOH}\): base fuerte, totalmente disociada. Es la disolución más básica (mayor pH).
            \item \(\ch{NH3}\): base débil, parcialmente disociada; su pH es básico, pero menor que el de una base fuerte de igual concentración.
            \item \(\ch{NaCl}\): sal de ácido fuerte y base fuerte, no se hidroliza; disolución neutra, \(pH = 7\).
            \item \(\ch{NH4Cl}\): sal de base débil y ácido fuerte; el \ch{NH4+} se hidroliza liberando \ch{H3O+} (\(\ch{NH4+ + H2O <=> NH3 + H3O+}\)), por lo que la disolución es ácida (menor pH de las cuatro).
        \end{itemize}

        Ordenando de menor a mayor pH:

        \begin{equation*}
            \ch{NH4Cl} < \ch{NaCl} < \ch{NH3} < \ch{NaOH}
        \end{equation*}

    \end{solu}
\end{exerciseEVAU}
\begin{exerciseEVAU}
    Se dispone de 50 mL de disolución acuosa 0,5 M de \ch{NH3}. a) Calcula el pH y el grado de disociación del amoniaco en la disolución. b) Escribe la reacción química que tiene lugar en la valoración de la disolución anterior con HCl 0,75 M. ¿Qué volumen del ácido se necesita para alcanzar el punto de equivalencia? c) ¿Qué volumen de amoniaco comercial del 24 \% (p/p) y densidad \densidad[gml]{0.91}  será necesario para preparar por dilución la disolución del enunciado?
    \datos{}: \(K_b(NH3)=\num{1.8e-5}\)); N = \num{14.0} u; H = \num{1.0} u.

    \begin{solu}

%        \begin{figure}[H]
%            \centering
%            \includegraphics[scale=.5]{images/imagesUD7EquilibrioAcidoBase/UD7_Problemas/EvAU_14.png}
%            %\caption{}
%            %\label{fig:}
%        \end{figure}

        \textbf{a) pH y grado de disociación}

        \begin{equation*}
            \ch{NH3 + H2O <=> NH4+ + OH-}
        \end{equation*}

        \begin{equation*}
            x = \sqrt{K_b\, C_0} = \sqrt{1,8\times10^{-5}\times0,5} = \num{3.0e-3} \ \si{\Molar}
        \end{equation*}

        \begin{equation*}
            \alpha = \dfrac{x}{C_0} = \dfrac{3,0\times10^{-3}}{0,5} = 0,006 \ (0,6\ \%)
        \end{equation*}

        \begin{equation*}
            pOH = -\log(3,0\times10^{-3}) = 2,52 \ \Rightarrow \ pH = 14-2,52 = 11,48
        \end{equation*}

        \textbf{b) Reacción de valoración y volumen de HCl}

        \begin{equation*}
            \ch{NH3 + HCl -> NH4Cl}
        \end{equation*}

        \begin{equation*}
            n(\ch{NH3}) = 0,050\times0,5 = 0,025 \ \text{mol} \ \Rightarrow \ V(\ch{HCl}) = \dfrac{0,025}{0,75} = 0,0333 \ \si{\litre} = \SI{33.3}{\milli\litre}
        \end{equation*}

        \textbf{c) Volumen de amoniaco comercial necesario}

        La masa de \ch{NH3} puro necesaria (\(M(\ch{NH3}) = \SI{17.0}{\gram\per\mole}\)) es:

        \begin{equation*}
            m(\ch{NH3}) = 0,025\times17,0 = \SI{0.425}{\gram}
        \end{equation*}

        Con una riqueza del 24\,\% (p/p):

        \begin{equation*}
            m_{disolución} = \dfrac{0,425}{0,24} = \SI{1.771}{\gram} \ \Rightarrow \ V = \dfrac{1,771}{0,91} = \SI{1.95}{\milli\litre}
        \end{equation*}

    \end{solu}
\end{exerciseEVAU}

\begin{exerciseEVAU}
    Escribe las fórmulas de las bases conjugadas de los siguientes ácidos: a) \ch{HCN} b) \ch{HCO3-} c) \ch{NH4+} d) \ch{HCl}

    \solutionCM{ a) \ch{CN-} ; b) \ch{CN\mch[2]} ; c) \ch{NH3} ; d) \ch{Cl-} }

    \begin{solu}

%        \begin{figure}[H]
%            \centering
%            \includegraphics[scale=.5]{images/imagesUD7EquilibrioAcidoBase/UD7_Problemas/EvAU_15.png}
%            %\caption{}
%            %\label{fig:}
%        \end{figure}

        La base conjugada de un ácido se obtiene eliminando un protón (\ch{H+}):

        \begin{itemize}
            \item a) \(\ch{HCN} \ \rightarrow \ \ch{CN-}\)
            \item b) \(\ch{HCO3-} \ \rightarrow \ \ch{CO3^{2-}}\)
            \item c) \(\ch{NH4+} \ \rightarrow \ \ch{NH3}\)
            \item d) \(\ch{HCl} \ \rightarrow \ \ch{Cl-}\)
        \end{itemize}

    \end{solu}
\end{exerciseEVAU}

\begin{exerciseEVAU}
    El ácido fluorhídrico está disociado al \num{0.5}\% en una disolución cuya concentración es \concentracion[m]{0.3}. Calcula: a) La constante de disociación del ácido. b) El pH de la disolución. c) La concentración molar de todos los iones presentes en la disolución.

    \solutionCM{a) \(\Ka = \num{7.54e-6}\); b) \(\pH = 2.82\); c) \(\conc{F-} = \conc{H3O+} = \concentracion[m]{1.5e-3}\) \conc{HF} = \concentracion[m]{0.2985} }

%\solutionCM{ a) \(K_a = \num{7.89e-4}\) ; b) pH = ; c) \con{HF} = \concentracion[m]{0.285} ; \conc{F-} = \conc{H3O+} = \concentracion[m]{0.015} }

    \begin{solu}

%        \begin{figure}[H]
%            \centering
%            \includegraphics[scale=.5]{images/imagesUD7EquilibrioAcidoBase/UD7_Problemas/EvAU_16.png}
%            %\caption{}
%            %\label{fig:}
%        \end{figure}

        Con \(\alpha = 0,005\) y \(C_0 = \SI{0.3}{\Molar}\), la concentración disociada es:

        \begin{equation*}
            x = \alpha\, C_0 = 0,005\times0,3 = \num{1.5e-3} \ \si{\Molar}
        \end{equation*}

        \textbf{a) Constante de disociación}

        \begin{equation*}
            \ch{HF + H2O <=> F- + H3O+}
        \end{equation*}

        \begin{equation*}
            K_a = \dfrac{x^2}{C_0-x} = \dfrac{(1,5\times10^{-3})^2}{0,3-0,0015} = \num{7.54e-6}
        \end{equation*}

        \textbf{b) pH de la disolución}

        \begin{equation*}
            pH = -\log(1,5\times10^{-3}) = 2,82
        \end{equation*}

        \textbf{c) Concentraciones de los iones}

        \begin{equation*}
            [\ch{F-}] = [\ch{H3O+}] = \num{1.5e-3} \ \si{\Molar} \qquad [\ch{HF}]_{eq} = 0,3-0,0015 = \SI{0.2985}{\Molar}
        \end{equation*}

    \end{solu}
\end{exerciseEVAU}
\begin{exerciseEVAU}
    El ácido fórmico o metanoico es un ácido monoprótico débil (\(\Ka = \num{1.8e-4}\) que se encuentra en la naturaleza en los aguijones de algunas especies de hormigas, de ahí su nombre. Se prepara una disolución \concentracion[m]{0.05}  de ácido fórmico. Calcula:
    \begin{enumerate}[label=\letraminuscula]
        \item El grado de disociación del ácido.
        \item El pH de la disolución.
        \item ¿Qué volumen de disolución de KOH 0,025 M será necesario para neutralizar 50 mL de la disolución de ácido fórmico? Escribe la reacción de neutralización.
        \item Razona en términos del principio de Le Chatelier qué le sucederá al grado de disociación del ácido si se incrementa el pH de la disolución, por ejemplo, añadiendo una base.
    \end{enumerate}


    \solutionCM{a) \(\alpha = 0.06\); b) \ph{2.52}; c) \(V_{\ch{KOH}}\) = \volumen[ml]{100} }

    \begin{solu}

%        \begin{figure}[H]
%            \centering
%            \includegraphics[scale=.5]{images/imagesUD7EquilibrioAcidoBase/UD7_Problemas/EvAU_17.png}
%            %\caption{}
%            %\label{fig:}
%        \end{figure}

        \begin{equation*}
            \ch{HCOOH + H2O <=> HCOO- + H3O+}
        \end{equation*}

        \textbf{a) Grado de disociación}

        \begin{equation*}
            x = \sqrt{K_a\, C_0} = \sqrt{1,8\times10^{-4}\times0,05} = \num{3.0e-3} \ \si{\Molar}
        \end{equation*}

        \begin{equation*}
            \alpha = \dfrac{x}{C_0} = \dfrac{3,0\times10^{-3}}{0,05} = 0,06 \ (6\ \%)
        \end{equation*}

        \textbf{b) pH de la disolución}

        \begin{equation*}
            pH = -\log(3,0\times10^{-3}) = 2,52
        \end{equation*}

        \textbf{c) Volumen de KOH necesario}

        \begin{equation*}
            \ch{HCOOH + KOH -> HCOOK + H2O}
        \end{equation*}

        \begin{equation*}
            n(\ch{HCOOH}) = 0,050\times0,05 = 0,0025 \ \text{mol} \ \Rightarrow \ V(\ch{KOH}) = \dfrac{0,0025}{0,025} = 0,1 \ \si{\litre} = \SI{100}{\milli\litre}
        \end{equation*}

        \textbf{d) Efecto de un aumento del pH (principio de Le Chatelier)}

        Si se aumenta el pH (por ejemplo, añadiendo una base), disminuye \([\ch{H3O+}]\). Como \([\ch{H3O+}]\) es uno de los productos del equilibrio de disociación del ácido, su disminución desplaza el equilibrio hacia la derecha para reponerlo parcialmente, es decir, \textbf{aumenta el grado de disociación} del ácido fórmico.

    \end{solu}
\end{exerciseEVAU}

\begin{exerciseEVAU}
    Justifica, escribiendo las ecuaciones químicas necesarias, por qué una disolución acuosa de \ch{NaClO} tiene \(\pH > 7\).

    \datos{\(\Ka (\ch{HClO}) = \num{2.9e-8}\)}


    \solutionCM{\ch{ClO- + H2O <=> HClO + OH-}}

    \begin{solu}

%        \begin{figure}[H]
%            \centering
%            \includegraphics[scale=.5]{images/imagesUD7EquilibrioAcidoBase/UD7_Problemas/EvAU_18.png}
%            %\caption{}
%            %\label{fig:}
%        \end{figure}

        El \ch{NaClO} es la sal de un ácido débil, el ácido hipocloroso \ch{HClO} (\(K_a = \num{2.9e-8}\)), y una base fuerte, \ch{NaOH}. El ion \ch{Na+} no se hidroliza, pero el ion hipoclorito, \ch{ClO-}, es la base conjugada de un ácido débil y sí se hidroliza en agua:

        \begin{equation*}
            \ch{ClO- + H2O <=> HClO + OH-}
        \end{equation*}

        Esta reacción produce un exceso de iones \ch{OH-} en la disolución. Como \(K_a(\ch{HClO})\) es muy pequeña, el equilibrio anterior (cuya constante es \(K_b = K_w/K_a\)) está bastante desplazado hacia la derecha, generando una concentración apreciable de \ch{OH-}. Por ello, una disolución acuosa de \ch{NaClO} tiene \(pH > 7\) (básico).

    \end{solu}
\end{exerciseEVAU}
\begin{exerciseEVAU}
    Ordena razonadamente de menor a mayor el pH de las disoluciones \concentracion[m]{0.1} de los siguientes compuestos: (a) \ch{NH4Cl}; (b) \ch{HCl}; (c) \ch{NaNO3}.


    \solutionCM{\ch{HCl} \(<\)  \ch{NH4Cl} \(<\) \ch{NaNO3}}

    \begin{solu}

%        \begin{figure}[H]
%            \centering
%            \includegraphics[scale=.5]{images/imagesUD7EquilibrioAcidoBase/UD7_Problemas/EvAU_19.png}
%            %\caption{}
%            %\label{fig:}
%        \end{figure}

        \begin{itemize}
            \item \(\ch{HCl}\): ácido fuerte, totalmente disociado. Es la disolución más ácida (menor pH).
            \item \(\ch{NH4Cl}\): sal de ácido fuerte y base débil; el ion \ch{NH4+} se hidroliza (\(\ch{NH4+ + H2O <=> NH3 + H3O+}\)), dando una disolución ácida, pero menos ácida que un ácido fuerte de igual concentración.
            \item \(\ch{NaNO3}\): sal de ácido fuerte y base fuerte, no se hidroliza; disolución neutra (\(pH = 7\)), la de mayor pH de las tres.
        \end{itemize}

        Ordenando de menor a mayor pH:

        \begin{equation*}
            \ch{HCl} < \ch{NH4Cl} < \ch{NaNO3}
        \end{equation*}

    \end{solu}
\end{exerciseEVAU}
\begin{exerciseEVAU}
    El pH medido en una botella de amoniaco doméstico es \num{11.97}. La etiqueta de dicha botella, indica que la densidad de la disolución de este amoniaco comercial es \densidad[gml]{0.97} . A partir de estos valores, calcula:
    \begin{enumerate}[label=\letraminuscula]
        \item La concentración del amoniaco en la disolución expresada en \% en volumen.
        \item El grado de ionización del mismo.
        \item ¿Cuál será el pOH que mediremos si diluimos 100 mL del amoniaco comercial con \volumen[ml]{150}de agua?
    \end{enumerate}
    \datos{}: \(k_b = \num{1.8e-5}\) ; masas atómicas: N = 14; H = 1; O = 16

    \solutionCM{a) \(11.23 \%\); b) \(\alpha = 0.0019\); \(\poh{2.52}\)}


    \begin{solu}

%        \begin{figure}[H]
%            \centering
%            \includegraphics[scale=.5]{images/imagesUD7EquilibrioAcidoBase/UD7_Problemas/EvAU_20.png}
%            %\caption{}
%            %\label{fig:}
%        \end{figure}

        A partir del pH obtenemos la concentración de \ch{OH-} en el amoniaco comercial:

        \begin{equation*}
            pOH = 14 - 11,97 = 2,03 \ \Rightarrow \ [\ch{OH-}] = 10^{-2,03} = \num{9.33e-3} \ \si{\Molar}
        \end{equation*}

        \textbf{a) Concentración del amoniaco comercial}

        Llamando \(x = [\ch{OH-}] = [\ch{NH4+}]\) y despejando la concentración inicial \(C_0\) a partir de \(K_b\):

        \begin{equation*}
            K_b = \dfrac{x^2}{C_0-x} \ \Rightarrow \ C_0 = x + \dfrac{x^2}{K_b} = 9,33\times10^{-3} + \dfrac{(9,33\times10^{-3})^2}{1,8\times10^{-5}} = \SI{4.85}{\Molar}
        \end{equation*}

        Expresando esta concentración como \% en masa, con \(M(\ch{NH3}) = \SI{17}{\gram\per\mole}\) y tomando 1 L de disolución como base de cálculo (masa de disolución \(= 1000 \times 0,97 = \SI{970}{\gram}\)):

        \begin{equation*}
            m(\ch{NH3}) = 4,85 \times 17 = \SI{82.4}{\gram} \ \Rightarrow \ \%(m/m) = \dfrac{82,4}{970}\times100 \approx 8,5\ \%
        \end{equation*}

        \textbf{b) Grado de ionización}

        \begin{equation*}
            \alpha = \dfrac{x}{C_0} = \dfrac{9,33\times10^{-3}}{4,85} = \num{1.92e-3} \ (0,19\ \%)
        \end{equation*}

        \textbf{c) pOH tras diluir 100 mL con 150 mL de agua}

        El factor de dilución es \((100+150)/100 = 2,5\), por lo que la nueva concentración inicial es:

        \begin{equation*}
            C_0' = \dfrac{4,85}{2,5} = \SI{1.94}{\Molar}
        \end{equation*}

        \begin{equation*}
            x' = \sqrt{K_b\, C_0'} = \sqrt{1,8\times10^{-5}\times1,94} = \num{5.91e-3} \ \si{\Molar}
        \end{equation*}

        \begin{equation*}
            pOH = -\log(5,91\times10^{-3}) = 2,23
        \end{equation*}

    \end{solu}
\end{exerciseEVAU}

\begin{exerciseEVAU}
    La Aspirina (ácido acetilsalicílico, \ch{C9H8O4}) es un ácido monoprótico débil, al cual podemos representar abreviadamente como HAsp, que da lugar al equilibrio de disociación:
    \begin{equation*}
        \ch{HAsp <=> H+ + Asp^{-}}
    \end{equation*}
    Para determinar experimentalmente su constante de disociación, un estudiante disolvió \masa[g]{2.00}  de aspirina en \volumen[ml]{600}  de agua encontrando que el pH de la disolución era \num{2.61}. Calcula:
    \begin{enumerate}[label=\letraminuscula]
        \item El valor de la constante de ionización \(K_a\)  de la aspirina.
        \item El grado de ionización de la misma en la disolución preparada por el estudiante.
        \item ¿Qué pH se habría medido si el alumno hubiese disuelto los 2.00 g de Aspirina en 1.00 L de agua?
    \end{enumerate}
    \datos{}: masas atómicas: C = 12, O= 16 y H=1

    \solutionCM{a) \(k_a = \num{3.74e-4}\); \(\alpha = 0.132\); c) \ph{2.73} }


    \begin{solu}

%        \begin{figure}[H]
%            \centering
%            \includegraphics[scale=.5]{images/imagesUD7EquilibrioAcidoBase/UD7_Problemas/EvAU_21.png}
%            %\caption{}
%            %\label{fig:}
%        \end{figure}

        La concentración inicial de aspirina es (\(M(\ch{C9H8O4}) = 9(12)+8(1)+4(16) = \SI{180}{\gram\per\mole}\)):

        \begin{equation*}
            n(\ch{HAsp}) = \dfrac{2,00}{180} = 0,01111 \ \text{mol} \ \Rightarrow \ C_0 = \dfrac{0,01111}{0,600} = \SI{0.01852}{\Molar}
        \end{equation*}

        \textbf{a) Constante de ionización \(K_a\)}

        \begin{equation*}
            [\ch{H+}] = 10^{-2,61} = \num{2.455e-3} \ \si{\Molar}
        \end{equation*}

        \begin{equation*}
            K_a = \dfrac{[\ch{H+}]^2}{C_0-[\ch{H+}]} = \dfrac{(2,455\times10^{-3})^2}{0,01852-0,002455} = \num{3.75e-4}
        \end{equation*}

        \textbf{b) Grado de ionización}

        \begin{equation*}
            \alpha = \dfrac{[\ch{H+}]}{C_0} = \dfrac{2,455\times10^{-3}}{0,01852} = 0,1326 \ (13,3\ \%)
        \end{equation*}

        \textbf{c) pH al disolver en 1,00 L}

        La nueva concentración inicial es:

        \begin{equation*}
            C_0' = \dfrac{0,01111}{1,00} = \SI{0.01111}{\Molar}
        \end{equation*}

        Resolviendo la ecuación de segundo grado \(K_a = \dfrac{x^2}{C_0'-x}\) con \(K_a = 3,75\times10^{-4}\):

        \begin{equation*}
            x^2 + 3,75\times10^{-4}\, x - 4,167\times10^{-6} = 0 \ \Rightarrow \ x = \num{1.86e-3} \ \si{\Molar}
        \end{equation*}

        \begin{equation*}
            pH = -\log(1,86\times10^{-3}) = 2,73
        \end{equation*}

    \end{solu}
\end{exerciseEVAU}
\begin{exerciseEVAU}
    Escribe los ácidos conjugados de las siguientes bases: (a) \ch{CN-} ; (b) \ch{HCO3-}; (c) \ch{N2H4}

    \solutionCM{a) \ch{HCN}; b) \ch{H2CO3} ; c) \ch{N2H5+}}

%
%  \begin{enumerate}[label=\letraminuscula]
%     \item ¿Qué volumen de disolución de \ch{KOH} \concentracion[m]{0,025} será necesario para neutralizar \volumen[ml]{50} de la disolución de ácido fórmico?.
%     \item Escribe la reacción de neutralización.
%     \item Razona en términos del principio de Le Chatelier qué le sucederá al grado de disociación del ácido si se incrementa el \(pH\) de la disolución, por ejemplo, añadiendo una base.
% \end{enumerate}

    \begin{solu}
%
%        \begin{figure}[H]
%            \centering
%            \includegraphics[scale=.5]{images/imagesUD7EquilibrioAcidoBase/UD7_Problemas/EvAU_22.png}
%            %\caption{}
%            %\label{fig:}
%        \end{figure}

        El ácido conjugado de una base se obtiene añadiendo un protón (\ch{H+}):

        \begin{itemize}
            \item a) \(\ch{CN-} \ \longrightarrow \ \ch{HCN}\)
            \item b) \(\ch{HCO3-} \ \longrightarrow \ \ch{H2CO3}\)
            \item c) \(\ch{N2H4} \ \longrightarrow \ \ch{N2H5+}\)
        \end{itemize}

    \end{solu}
\end{exerciseEVAU}

